% =========================
% Sections 6–7 (LaTeX-ready)
% =========================
% Minimal preamble suggestions:
% \usepackage{amsmath,amssymb,amsthm}
% \usepackage{algorithm}
% \usepackage{algpseudocode} % or algorithmicx
% (Optional for nicer boxes: \usepackage[most]{tcolorbox}

% ---------- Theorem environments ----------
\newtheorem{theorem}{Theorem}
\newtheorem{corollary}{Corollary}
\newtheorem{lemma}{Lemma}
\theoremstyle{definition}
\newtheorem{definition}{Definition}
\newtheorem{assumption}{Assumption}

% ---------- A simple "Assumptions box" (no extra packages) ----------
\newcommand{\StandingAssumptions}{%
\begin{center}
\fbox{%
\begin{minipage}{0.97\linewidth}
\textbf{Standing Assumptions (A–D).}
\begin{enumerate}
\item[\textbf{A.}]\textbf{(Space + norm).} $(\mathcal X,\|\cdot\|)$ is a Banach space over $\mathbb K\in\{\mathbb R,\mathbb C\}$.
Operator norms are induced: $\|A\|=\sup_{\|x\|=1}\|Ax\|$ for $A\in\mathcal B(\mathcal X)$.
\item[\textbf{B.}]\textbf{(Forward invariance).} There exists a nonempty closed set $\Omega\subseteq\mathcal X$ such that for all $t$,
\[
x\in\Omega \implies \Phi_t(x)\in\Omega,\qquad \Phi_t(x)=\Xi(t)x+\Lambda_m(t)\,T(x).
\]
\item[\textbf{C.}]\textbf{(Bounded sector operator).} For each $t$,
\[
\Xi(t)=\sum_{p\in P_N} a_p(t)\,U_p(t)\in\mathcal B(\mathcal X)
\]
is well-defined and bounded, with the computable bound
\[
\|\Xi(t)\|\le \sum_{p\in P_N} |a_p(t)|\,\|U_p(t)\|.
\]
\item[\textbf{D.}]\textbf{(Nonlinearity regularity).} $T:\Omega\to\mathcal X$ is $L_T$-Lipschitz on $\Omega$
(or admits a known upper bound $L_T$ on $\Omega$).
\end{enumerate}
\vspace{0.25em}
\textbf{Note (Prime indexing).} Prime indexing is a deterministic sparse schedule; all correctness results hold for any finite index set.
Primes are used only as an engineering choice for reproducible sparsity.
\end{minipage}}
\end{center}
}

% ============================================================
\section{Stability Results: Uniform Incremental Contraction}
% ============================================================

\StandingAssumptions

\begin{definition}[DRMM update map]
For each time step $t$, define
\[
\Phi_t(x):=\Xi(t)x+\Lambda_m(t)\,T(x),\qquad x\in\Omega.
\]
Trajectories satisfy $X_{t+1}=\Phi_t(X_t)$.
\end{definition}

\begin{theorem}[Uniform incremental contraction]\label{thm:incremental_contraction}
Assume Standing Assumptions (A--D). Suppose there exist $\varepsilon\in(0,1)$ and $c\in[0,\varepsilon)$ such that for all $t$,
\[
\|\Xi(t)\|\le 1-\varepsilon,\qquad \Lambda_m(t)\,L_T\le c.
\]
Then for all $x,y\in\Omega$,
\[
\|\Phi_t(x)-\Phi_t(y)\|\le \rho\,\|x-y\|,\qquad \rho:=1-\varepsilon+c\in(0,1),
\]
and for any two trajectories $(X_t)$ and $(Y_t)$ driven by the same time-varying maps,
\[
\|X_t-Y_t\|\le \rho^t\,\|X_0-Y_0\|,\qquad t\ge 0.
\]
\end{theorem}

\begin{proof}
For $x,y\in\Omega$,
\[
\|\Phi_t(x)-\Phi_t(y)\|
=\|\Xi(t)(x-y)+\Lambda_m(t)(T(x)-T(y))\|
\le \|\Xi(t)\|\,\|x-y\|+\Lambda_m(t)\,\|T(x)-T(y)\|.
\]
By Lipschitzness of $T$ on $\Omega$, $\|T(x)-T(y)\|\le L_T\|x-y\|$, hence
\[
\|\Phi_t(x)-\Phi_t(y)\|\le(\|\Xi(t)\|+\Lambda_m(t)L_T)\|x-y\|\le (1-\varepsilon+c)\|x-y\|=\rho\|x-y\|.
\]
Iterating yields $\|X_t-Y_t\|\le\rho^t\|X_0-Y_0\|$.
\end{proof}

\begin{corollary}[Time-invariant fixed point]\label{cor:fixed_point}
Assume Standing Assumptions (A--D) and the hypotheses of Theorem~\ref{thm:incremental_contraction}.
If additionally $\Xi(t)\equiv \Xi$ and $\Lambda_m(t)\equiv\Lambda_m$ are constant in $t$, then
$\Phi(x):=\Xi x+\Lambda_m T(x)$ is a contraction on $\Omega$, hence has a unique fixed point $x^\star\in\Omega$,
and every trajectory $X_{t+1}=\Phi(X_t)$ with $X_0\in\Omega$ converges to $x^\star$.
\end{corollary}

\begin{corollary}[Optional: convergent/slowly varying maps (tracking)]\label{cor:tracking}
Assume the hypotheses of Theorem~\ref{thm:incremental_contraction}. If there exists a limit map $\Phi_\infty:\Omega\to\Omega$
such that $\sup_{x\in\Omega}\|\Phi_t(x)-\Phi_\infty(x)\|\to 0$ as $t\to\infty$, and all $\Phi_t$ share the same contraction modulus
$\rho<1$, then the DRMM trajectory $(X_t)$ tracks the unique fixed point of $\Phi_\infty$ (standard contracting time-varying system tracking).
\end{corollary}

% ------------------------------------------------------------
\subsection{A referee-friendly forward-invariance sufficient condition}
% ------------------------------------------------------------

\begin{lemma}[Ball invariance via a growth bound]\label{lem:ball_invariance}
Let $\Omega=B_R:=\{x\in\mathcal X:\|x\|\le R\}$. Suppose that for all $\|x\|\le R$ and all $t$,
\[
\|\Xi(t)x\|\le \|\Xi(t)\|\,R,\qquad \|T(x)\|\le \beta+L_T\|x\|
\]
for some $\beta\ge 0$ and $L_T\ge 0$. If
\[
\|\Xi(t)\|R+\Lambda_m(t)\big(\beta+L_T R\big)\le R\qquad\text{for all }t,
\]
then $B_R$ is forward-invariant: $x\in B_R\implies \Phi_t(x)\in B_R$ for all $t$.
\end{lemma}

\begin{proof}
If $\|x\|\le R$, then
\[
\|\Phi_t(x)\|\le \|\Xi(t)x\|+\Lambda_m(t)\|T(x)\|
\le \|\Xi(t)\|R+\Lambda_m(t)\big(\beta+L_T R\big)\le R.
\]
\end{proof}

% ============================================================
\section{Certified Spectral Control (CSC): Per-step Certificates and Gain Clamping}
% ============================================================

\begin{definition}[Structural vs.\ analytic $\Lambda$ (spec-level reconciliation)]
Let $\widehat{\Lambda}(t)$ denote a \emph{structural} multiplicity/PGF functional (descriptive; not used as a gain).
Let $\Lambda_m(t)\in\mathbb R_{\ge 0}$ denote the \emph{analytic certification gain} chosen (or clamped) by CSC to enforce contraction.
\end{definition}

\begin{definition}[Per-step certificate quantities]
Given an upper bound $\widehat{\|\Xi(t)\|}$ and a Lipschitz bound $\widehat{L_T(t)}$ (global or local), define the
\emph{certified upper bound} on the contraction modulus:
\[
\widehat{\rho}(t):=\widehat{\|\Xi(t)\|}+\Lambda_m(t)\,\widehat{L_T(t)}.
\]
Define also
\[
\widehat{\varepsilon}(t):=1-\widehat{\|\Xi(t)\|},\qquad \widehat{c}(t):=\Lambda_m(t)\,\widehat{L_T(t)}.
\]
A CSC report is $\mathrm{CSC}(t):=(\widehat{\varepsilon}(t),\widehat{c}(t),\widehat{\rho}(t))$ with pass/fail given by $\widehat{\rho}(t)<1$.
\end{definition}

\begin{algorithm}[t]
\caption{CSC (per-step certificate + gain clamping; global/local/hybrid)}
\label{alg:csc}
\begin{algorithmic}[1]
\Require Prime index set $P_N$; sector operators $U_p(t)$; weights $a_p(t)$; proposal $\widehat{\Lambda}(t)$; safety margin $\delta>0$;
gain cap $\Lambda_{\max}>0$; global Lipschitz bound $L_T^{\mathrm{glob}}$ (optional); local estimator for $\widehat{L_T(t)}$ (optional).
\Ensure Certified gain $\Lambda_m(t)$ and certificate tuple $(\widehat{\varepsilon}(t),\widehat{c}(t),\widehat{\rho}(t))$.
\vspace{0.25em}

\State \textbf{(Linear bound)} Compute an upper bound
\[
\widehat{\|\Xi(t)\|}\ \leftarrow\ \sum_{p\in P_N}|a_p(t)|\,\|U_p(t)\|
\quad\text{(or any tighter computable bound).}
\]
\State Set $\widehat{\varepsilon}(t)\leftarrow 1-\widehat{\|\Xi(t)\|}$.
\State \textbf{(Lipschitz bound)} If a reliable local estimate is available, set $\widehat{L_T(t)}\leftarrow$ local estimate; else set $\widehat{L_T(t)}\leftarrow L_T^{\mathrm{glob}}$.
\State \textbf{(Clamp target)} Set
\[
\Lambda_{\mathrm{target}}(t)\leftarrow \max\left\{0,\ \frac{\widehat{\varepsilon}(t)-\delta}{\widehat{L_T(t)}}\right\}.
\]
\State \textbf{(Clamp)} Set
\[
\Lambda_m(t)\leftarrow \min\{\Lambda_{\max},\ \Lambda_{\mathrm{target}}(t),\ \widehat{\Lambda}(t)\}.
\]
\State \textbf{(Certificate)} Set $\widehat{c}(t)\leftarrow \Lambda_m(t)\widehat{L_T(t)}$ and $\widehat{\rho}(t)\leftarrow \widehat{\|\Xi(t)\|}+\widehat{c}(t)$.
\State \textbf{(Audit log)} Record $\widehat{\|\Xi(t)\|}$ (and bound method), $\widehat{L_T(t)}$ (and method), $\widehat{\Lambda}(t)$, $\Lambda_m(t)$, $\widehat{\varepsilon}(t)$, $\widehat{c}(t)$, $\widehat{\rho}(t)$, and pass/fail $(\widehat{\rho}(t)<1)$.
\State \Return $(\Lambda_m(t),\widehat{\varepsilon}(t),\widehat{c}(t),\widehat{\rho}(t))$.
\end{algorithmic}
\end{algorithm}

\paragraph{Logging contract (auditable punchline).}
For each step $t$, CSC must log:
(i) the bound used for $\widehat{\|\Xi(t)\|}$ (e.g.\ $\sum_p |a_p(t)|\|U_p(t)\|$),
(ii) the bound/estimate used for $\widehat{L_T(t)}$ (global vs local and method),
(iii) the proposal $\widehat{\Lambda}(t)$ and the certified gain $\Lambda_m(t)$,
(iv) the certificate tuple $(\widehat{\varepsilon}(t),\widehat{c}(t),\widehat{\rho}(t))$ and pass/fail.
